\section{การตรวจสอบสภาพถนนด้วยเซ็นเซอร์สมาร์ตโฟน}

\subsection{กระบวนการและเซ็นเซอร์ที่ใช้}

Sattar et al. (2018) สรุปกระบวนการตรวจสอบสภาพถนนด้วยเซ็นเซอร์สมาร์ตโฟนไว้เป็น 5 ขั้นตอนหลัก ได้แก่ การเก็บข้อมูลจากเซ็นเซอร์ การเตรียมข้อมูลล่วงหน้า การสกัดคุณลักษณะ การประมวลผลหลัง และการประเมินประสิทธิภาพ ขณะที่ Alqaydi et al. (2024) ชี้ให้เห็นว่าแนวทางสมาร์ตโฟนร่วมกับการประมวลผลข้อมูลและแบบจำลองการเรียนรู้ของเครื่องสามารถลดต้นทุนและความซับซ้อนของการสำรวจถนนด้วยวิธีดั้งเดิมได้

เซ็นเซอร์หลักที่ใช้ในงานวิจัยด้านนี้ประกอบด้วยมาตรวัดความเร่ง ซึ่งวัดความเร่งใน 3 แกน มาตรวัดการหมุน ซึ่งวัดอัตราการหมุนเชิงมุมใน 3 แกน มาตรวัดสนามแม่เหล็ก ซึ่งช่วยระบุทิศทางของอุปกรณ์ และระบบระบุตำแหน่งบนพื้นโลก ซึ่งให้ข้อมูลตำแหน่ง ความเร็ว และทิศทางการเคลื่อนที่ ปัจจัยสำคัญที่ส่งผลต่อคุณภาพของสัญญาณคือ ตำแหน่งและทิศทางการติดตั้งสมาร์ตโฟนภายในรถ เนื่องจากสัญญาณจากเหตุการณ์เดียวกันอาจปรากฏในแกนที่แตกต่างกันตามทิศทางของอุปกรณ์ ทำให้แบบจำลองที่ฝึกจากตำแหน่งการติดตั้งหนึ่งอาจมีประสิทธิภาพลดลงเมื่อใช้กับตำแหน่งการติดตั้งที่แตกต่างกัน งานวิจัยก่อนหน้าจึงให้ความสำคัญกับการกรองสัญญาณ การปรับแกนอ้างอิง และการแปลงพิกัดจากกรอบอ้างอิงของสมาร์ตโฟนไปสู่กรอบอ้างอิงของยานพาหนะก่อนนำข้อมูลเข้าสู่แบบจำลอง (Sattar et al., 2018; Alqaydi et al., 2024)

\subsection{วิธีการจำแนกสภาพถนนในงานวิจัยที่ผ่านมา}

งานวิจัยในระยะแรกส่วนใหญ่ใช้วิธีการแบบอาศัยค่าเกณฑ์ ซึ่งกำหนดเงื่อนไขการตรวจจับโดยตรงจากสัญญาณมาตรวัดความเร่ง โดยเฉพาะสัญญาณความเร่งในแนวดิ่ง Mednis et al. (2011) เสนอกฎสำหรับตรวจจับหลุมบ่อบนถนนในประเทศลัตเวียหลายรูปแบบ ได้แก่ Z-THRESH, Z-DIFF และ G-ZERO โดยพบว่าวิธี Z-THRESH ให้ความถูกต้องในการตรวจจับสูงถึงร้อยละ 90 ภายใต้สภาวะการทดลองที่ควบคุมได้ ขณะที่ Perttunen et al. (2011) ประยุกต์ใช้พลังงานสเปกตรัมจากการแปลงฟูเรียร์แบบเร็วร่วมกับค่าเกณฑ์ เพื่อประเมินความขรุขระของถนนในประเทศฟินแลนด์ อย่างไรก็ตาม ข้อจำกัดร่วมของวิธีการกลุ่มนี้คือ ค่าเกณฑ์ที่เหมาะสมมักขึ้นอยู่กับประเภทยานพาหนะ ความเร็ว และตำแหน่งการติดตั้งสมาร์ตโฟน ทำให้ค่าเกณฑ์ที่ปรับจากชุดข้อมูลหนึ่งอาจไม่สามารถนำไปใช้กับสภาวะแวดล้อมอื่นได้โดยตรง (Sattar et al., 2018; Alqaydi et al., 2024)

แนวทางถัดมาประยุกต์ใช้การเรียนรู้ของเครื่องแบบดั้งเดิมร่วมกับการสกัดคุณลักษณะทางสถิติจากสัญญาณ Gonzalez et al. (2017) เสนอการแทนข้อมูลแบบ Bag-of-Words เพื่อแปลงสัญญาณมาตรวัดความเร่งแบบอนุกรมเวลาให้อยู่ในรูปเวกเตอร์คุณลักษณะ ก่อนนำเข้าสู่ตัวจำแนกหลายรูปแบบ โดยเก็บข้อมูลจากยานพาหนะ 12 คัน และครอบคลุมตำแหน่งการวางสมาร์ตโฟนหลายรูปแบบในเมืองชิวาวา ประเทศเม็กซิโก Basavaraju et al. (2020) เปรียบเทียบเครื่องเวกเตอร์สนับสนุน (Support Vector Machine: SVM) ต้นไม้ตัดสินใจ และโครงข่ายประสาทเทียม สำหรับจำแนกสภาพถนน 3 ประเภท ได้แก่ ถนนเรียบ หลุมบ่อ และรอยแตกร้าว บนข้อมูลภาคสนามจากสหรัฐอเมริกา โดยพบว่าโครงข่ายประสาทเทียมให้ความถูกต้องประมาณร้อยละ 92 วิธีการกลุ่มนี้มีข้อดีด้านความเรียบง่ายและการตีความผลเมื่อเทียบกับวิธีการเรียนรู้เชิงลึก แต่ยังต้องอาศัยการออกแบบคุณลักษณะด้วยมือ ซึ่งอาจส่งผลต่อความสามารถในการนำไปใช้กับสภาพแวดล้อมที่แตกต่างจากชุดข้อมูลเดิม

งานวิจัยกลุ่มล่าสุดเริ่มประยุกต์ใช้การเรียนรู้เชิงลึก โดยรับสัญญาณจากหน่วยวัดความเฉื่อยเป็นข้อมูลนำเข้าโดยตรง Chen et al. (2023) ใช้สถาปัตยกรรมโครงข่ายประสาทเทียมแบบคอนโวลูชัน (Convolutional Neural Network: CNN) ร่วมกับโครงข่ายหน่วยความจำระยะสั้นแบบยาว (Long Short-Term Memory: LSTM) เพื่อจำแนกพื้นผิวถนนของยานพาหนะไฟฟ้าในประเทศจีน งานในกลุ่มนี้แสดงให้เห็นว่าการเรียนรู้เชิงลึกสามารถเรียนรู้รูปแบบจากสัญญาณดิบได้โดยลดการพึ่งพาการออกแบบคุณลักษณะด้วยมือ อย่างไรก็ตาม งานวิจัยแต่ละชิ้นยังมีความแตกต่างด้านบริบทพื้นที่ศึกษา ประเภทยานพาหนะ อุปกรณ์ที่ใช้เก็บข้อมูล ตำแหน่งการติดตั้งสมาร์ตโฟน และระเบียบวิธีการกำหนดป้ายกำกับ อีกทั้งหลายงานไม่ได้เปรียบเทียบแบบจำลองหลายรูปแบบภายใต้ชุดข้อมูลและกระบวนการประเมินเดียวกัน จึงยังยากต่อการสรุปว่าแบบจำลองใดเหมาะสมที่สุดสำหรับการจำแนกสภาพถนนในสภาพแวดล้อมจริง งานวิจัยที่เกี่ยวข้องสรุปไว้ในตารางที่ 2.2

\begingroup
\fontsize{11pt}{13pt}\selectfont
\setlength{\tabcolsep}{2pt}
\renewcommand{\arraystretch}{1.15}
\setlength{\LTleft}{0pt}
\setlength{\LTright}{0pt}
\begin{longtable}{>{\raggedright\arraybackslash}p{0.14\textwidth}>{\raggedright\arraybackslash}p{0.18\textwidth}>{\raggedright\arraybackslash}p{0.17\textwidth}>{\raggedright\arraybackslash}p{0.15\textwidth}>{\raggedright\arraybackslash}p{0.14\textwidth}>{\raggedright\arraybackslash}p{0.13\textwidth}}
\caption{สรุปงานวิจัยด้านการตรวจสอบสภาพถนนด้วยสมาร์ตโฟน}\label{tab:smartphone-road-research}\\
\toprule
\textbf{ผู้วิจัย (ปี)} & \textbf{Method} & \textbf{Dataset} & \textbf{Classes} & \textbf{Accuracy/F1} & \textbf{Dataset type} \\
\midrule
\endfirsthead
\multicolumn{6}{l}{\textit{ตารางที่ \thetable\ (ต่อ)}}\\
\toprule
\textbf{ผู้วิจัย (ปี)} & \textbf{Method} & \textbf{Dataset} & \textbf{Classes} & \textbf{Accuracy/F1} & \textbf{Dataset type} \\
\midrule
\endhead
\bottomrule
\endfoot
Mednis et al. (2011) & Threshold-based (Z-THRESH, Z-DIFF, G-ZERO) & Private (Latvia) & 2 (Normal, Pothole) & TPR 90\% & Real-world field \\
Perttunen et al. (2011) & FFT Spectral Energy + Threshold & Private (Finland) & Road roughness levels & n.r. & Real-world field \\
Gonzalez et al. (2017) & Bag-of-Words (BoW) + ML classifiers & Public (Mexico, Chihuahua) & 5 (Pothole, Metal Bump, Asphalt Bump, Worn-out, Regular) & n.r. & Real-world field \\
Silva et al. (2018) & System evaluation (real-world deployment analysis) & Private (Portugal, controlled + real-world) & Road anomalies & Performance degradation documented: controlled → real-world & Real-world field \\
Basavaraju et al. (2020) & SVM, Decision Tree, Neural Network & Private (USA) & 3 (Smooth, Pothole, Crack) & Acc ≈ 92\% (NN) & Real-world field \\
Chibani et al. (2022) & DSW vs SSW + ML classifiers (SVM, DT, GB, RF, NB, NN, KNN) & Pothole Lab + Carsim (benchmark/generated) & Road anomalies & DSW outperforms SSW; KNN F1: 0.718→0.862, RF F1: 0.663→0.832 & Benchmark / Generated \\
Chen et al. (2023) & CNN-LSTM & Private (China, EBV) & ≥ 2 surface types & n.r. & Real-world field \\
\end{longtable}
\endgroup
